%! Multiplying \& Dividing Radicals; Rationalizing — Unit 6, Lesson 6.3 — Lesson Plan
\documentclass[10pt]{article}
\usepackage{algebra2-article}
\usepackage{algebra2-boxes}
\usepackage{graphicx}   % -article does NOT load graphicx; needed for thumbnails/screenshots
\graphicspath{{images/}}
\newcommand{\TallMath}[1]{$\displaystyle #1\rule[-1.4em]{0pt}{3.2em}$}
\newcommand{\CourseName}{Algebra 2: Shepherd}
\newcommand{\SchoolYear}{2026--2027}
\newcommand{\MeetingLength}{55 minutes}
\newcommand{\UnitNumberName}{Unit 6: Radical Functions \quad}
\newcommand{\LessonNumberName}{Lesson 6.3: Multiplying \& Dividing Radicals; Rationalizing}

\begin{document}
\begin{center}
    {\Large\bfseries \CourseName: \SchoolYear \\
    {\small\bfseries \UnitNumberName \LessonNumberName}}
\end{center}

\begin{tcolorbox}[colback=forestbg, colframe=forest, boxrule=0.9pt, arc=2mm,
    left=2mm, right=2mm, top=1.5mm, bottom=1.5mm]
    \textbf{Primary Objective:} Students can run the product and quotient properties of radicals
    \emph{forwards} --- multiplying radicals, distributing across a sum, and multiplying two
    binomials --- and can simplify the result, recognizing that a product frequently contains a
    perfect power that neither factor had. They can state and use the single mechanical fact the
    lesson rests on, $\sqrt{a}\cdot\sqrt{a}=a$; can identify a \emph{conjugate} and explain why a
    conjugate pair always multiplies to a \emph{rational} number, while the same binomial squared
    does not; and can \emph{rationalize} a denominator --- one radical, a higher index, or a
    binomial --- while explaining why multiplying by a form of $1$ changes an expression's form but
    never its value.
    \\[2pt]
    \textbf{Standards (2023 VA SOL):} \textbf{A2.EO.2b} --- add, subtract, multiply, and divide
    radical expressions, including rationalizing denominators --- is the governing standard, and
    this lesson carries its \emph{multiplicative} half (Lesson 6.2 carried the additive half; that
    split is why A2.EO.2b spans the pair). \textbf{A2.EO.2a} is used continuously, since every
    answer must be returned to simplest radical form, and today closes simplest-form condition 3,
    which 6.2 deliberately left open. \textbf{A2.EO.2c} (radical $\leftrightarrow$ rational
    exponent) reappears in the mismatched-index case $\sqrt2\cdot\sqrt[3]2$ and in the Activity's
    and homework's ``two roads'' items, where rationalizing is shown to be nothing but converting
    $a^{-1/n}$ into a positive exponent. The conjugate work is a deliberate structural callback to
    \textbf{A2.EO.4c} (Lesson 3.4), where $(a+bi)(a-bi)$ served the identical purpose in a
    different number system.
\end{tcolorbox}

\renewcommand{\arraystretch}{1.4}
\begin{skillbox}[Priority Ideas \& Skills]{goldbox}
\begin{tabularx}{\linewidth}{@{}X|X@{}}
\textbf{Priority Ideas \& Skills} & \textbf{Key Understandings (the ``why'')} \\[2pt]
\hline
\vspace{-6pt}
\begin{itemize}[leftmargin=1.1em, nosep, after=\vspace{2pt}]
  \item Multiply radicals: outsides times outsides, insides times insides, \textbf{then} simplify.
  \item \textbf{Distribute} a radical over a sum and \textbf{FOIL} two radical binomials.
  \item Square a radical binomial \emph{with} its middle term.
  \item Write the \textbf{conjugate} of $a+\sqrt b$ and predict that the product is rational.
  \item Divide radicals with the quotient property, run forwards.
  \item \textbf{Rationalize} a monomial denominator --- including index $>2$.
  \item \textbf{Rationalize} a binomial denominator with the conjugate.
  \item Justify rationalizing as multiplication by a form of $1$.
\end{itemize}
&
\vspace{-6pt}
\begin{itemize}[leftmargin=1.1em, nosep, after=\vspace{2pt}]
  \item Nothing here is a new rule. It is yesterday's property read right-to-left, plus the
        distributive property students have used since Algebra 1.
  \item \textbf{$\sqrt a\cdot\sqrt a=a$}: a square root times itself erases the radical. Every
        other idea today is a consequence of this one line.
  \item A \emph{product} often simplifies when neither factor could --- so ``finish it'' is a habit,
        not a step.
  \item Conjugates work because the middle terms \emph{cancel} and the squares \emph{clear}. Change
        one sign and the effect vanishes.
  \item Rationalizing is legal because $\frac{\sqrt b}{\sqrt b}=1$. The form changes; the value
        cannot. Without that, it is a ritual.
  \item The reason the rule exists at all is arithmetic: dividing \emph{by} an irrational number is
        painful; dividing by a whole number is not.
\end{itemize}
\end{tabularx}
\end{skillbox}

\begin{skillbox}[{Vocabulary, Concepts, \& Theorems}]{sky}
\renewcommand{\arraystretch}{1.3}
\begin{tabularx}{\linewidth}{@{}l X@{}}
\textbf{Product Property (forwards)} & \TallMath{\sqrt[n]{a}\cdot\sqrt[n]{b}=\sqrt[n]{ab}} --- same statement as Lesson 6.2, read the other direction. Indices must match. \\
\textbf{Quotient Property (forwards)} & \TallMath{\dfrac{\sqrt[n]{a}}{\sqrt[n]{b}}=\sqrt[n]{\dfrac{a}{b}}}, $b\ne0$ --- two radicals with one index pull together under one radical. \\
\textbf{The eraser} & $\sqrt{a}\cdot\sqrt{a}=a$ and $\left(\sqrt a\right)^{2}=a$. A cube root needs \emph{three} copies: the index still sets the group size. \\
\textbf{Conjugate} & The same binomial with the middle sign reversed: $a+\sqrt b$ and $a-\sqrt b$. \\
\textbf{Conjugate product} & $\left(a+\sqrt b\right)\left(a-\sqrt b\right)=a^{2}-b$ and $\left(\sqrt a+\sqrt b\right)\left(\sqrt a-\sqrt b\right)=a-b$ --- \emph{always rational}. \\
\textbf{Rationalizing} & Clearing a radical out of a denominator by multiplying the numerator \emph{and} denominator by the same expression. \\
\textbf{Form of $1$} & $\frac{\sqrt b}{\sqrt b}=1$; multiplying by it changes an expression's \emph{form}, never its \emph{value}. The justification for the whole method. \\
\textbf{Simplest radical form} & Now complete: (1) no perfect $n$th power factor; (2) no fraction under a radical; (3) no radical in a denominator. \\
\end{tabularx}
\end{skillbox}

\begin{fixedskillbox}[Activate Prior Knowledge \& Spiral Review]{forestbg}
    The Warm-Up rehearses three prerequisites, and each is a third of today's lesson in disguise.
    \textbf{(1)} Distributing and multiplying binomials --- $3x(2x+5)$, $(x+4)(x-4)$, $(x+4)^{2}$.
    The follow-up question about $(x+4)(x-4)$ is doing the most work on the page: get the words
    \emph{``same terms, opposite middle sign''} said aloud now, because that sentence \emph{is} the
    definition of a conjugate, arriving in Section 3. And $(x+4)^{2}$ pre-kills the middle-term
    error before it can happen --- when a student later writes
    $\left(\sqrt5+\sqrt2\right)^{2}=7$, you can point at work they did themselves this period.
    \textbf{(2)} $\sqrt{12}$, $\sqrt{50}$ from yesterday, then $\sqrt7\cdot\sqrt7$ and
    $\left(\sqrt{11}\right)^{2}$ --- the entire engine of the lesson. Make sure the word
    \emph{erases} (or ``undoes,'' ``cancels the radical'') is spoken. \textbf{(3)} Equivalent
    fractions: $\frac34=\frac{15}{20}$, and $\frac55=1$. This is the honesty check on rationalizing;
    students who hesitate on the last blank are the ones who will suspect Section 5 of cheating.
\end{fixedskillbox}

\begin{skillbox}[Hook]{forestbg}
    Put two expressions on the board and ask the class to decide which student is right.
    \begin{center}
    Student A: \; $\dfrac{1}{\sqrt2}$ \qquad\qquad Student B: \; $\dfrac{\sqrt2}{2}$
    \end{center}
    Let them argue for a minute, then hand out calculators: both are $0.7071067\ldots$ Neither
    student is wrong. \textbf{So ask the honest question: ``If both are correct, why does every
    textbook insist on the second one?''} Students will guess ``it looks nicer'' or ``it's the
    rule.'' Give them the real answer instead, because it is the one that makes the rest of the
    period feel non-arbitrary: \emph{``It is 1850. You have no calculator, and you must actually
    compute one of these by hand. Which long division would you rather do ---
    $1\div1.414214$, or $1.414214\div2$?''} Every hand agrees. That is the whole reason simplest
    form has a third condition: a radical on \emph{top} is harmless, a radical on the \emph{bottom}
    is a nuisance, and the convention long predates the calculator that made it cosmetic.
    \\[3pt]
    Then turn to the mechanism: \textbf{how did Student B get there?} They multiplied top and
    bottom by $\sqrt2$, and $\sqrt2\cdot\sqrt2=2$ --- the radical simply left. Point back at
    Warm-Up item 2 and Warm-Up item 3: they have already established both halves of that move
    (a root times itself erases; multiplying by $\frac{n}{n}$ is safe). Today is those two facts,
    applied everywhere.
\end{skillbox}

\begin{skillbox}[Lesson]{forestbg}
\begin{multicols}{2}
\textbf{1. Multiplying (about 8 min).} Do not present it as new: write
$\sqrt[n]{a}\cdot\sqrt[n]{b}=\sqrt[n]{ab}$ and say out loud that it is yesterday's property read
right-to-left. The teaching point is the \emph{last} step, not the first: choose examples where
neither factor simplifies alone but the product does ($\sqrt6\cdot\sqrt{10}=\sqrt{60}=2\sqrt{15}$),
so that ``finish it'' is discovered rather than nagged. Then isolate $\sqrt a\cdot\sqrt a=a$ and
name it --- \emph{the eraser} --- and note the cube-root version needs three copies. Mention
mismatched indices once ($\sqrt2\cdot\sqrt[3]2=2^{5/6}$, Lesson 6.1) and move on.

\textbf{2. Distributing \& FOIL (about 6 min).} Fast, because it is Algebra 1 with radicals
substituted for variables. Spend the time only on $\left(\sqrt7+\sqrt2\right)^{2}$, and connect it
to Warm-Up 1(c) explicitly.

\columnbreak

\textbf{3. Conjugates (about 10 min --- the center of the lesson).} \emph{Do not state the rule
first.} Work $\left(4+\sqrt3\right)\left(4-\sqrt3\right)$ term by term on the board, let the class
watch $-4\sqrt3$ and $+4\sqrt3$ cancel, and then ask: ``what kind of number is $13$?'' The word
\emph{rational} should come from a student. Immediately put $\left(4+\sqrt3\right)^{2}=19+8\sqrt3$
beside it --- same two numbers, radical survives. That contrast is the argument. Then name the
term, generalize to $a^{2}-b$, and spend ninety seconds on the Lesson 3.4 callback:
$(a+bi)(a-bi)=a^{2}+b^{2}$ did the same job in a different number system, and the sign difference is
entirely $i^{2}=-1$ versus $\left(\sqrt b\right)^{2}=+b$.

\textbf{4--5. Dividing, then rationalizing one radical (about 8 min).} Division is quick; land on
$\sqrt3/\sqrt5$ as the one that will not come out, which reopens condition 3 by name. Then the
move, with the justification \emph{first}: $\frac{\sqrt b}{\sqrt b}=1$. Include the higher-index
case $\frac{1}{\sqrt[3]2}$ --- it needs $\sqrt[3]4$, not $\sqrt[3]2$, and the reason is yesterday's
``the index sets the group size,'' now read forwards.

\textbf{6. Binomial denominators (about 12 min).} Ask first: \emph{what could we possibly multiply
$3-\sqrt5$ by?} Let someone propose $\sqrt5$ and let it fail on the board. Then the conjugate,
which they now have a reason to expect. Finish with the reduce-at-the-end habit and the three
errors.
\end{multicols}
\end{skillbox}

\begin{skillbox}[Active Monitoring]{redbox}
Circulate during Guided Practice and the tiers. Watch for and correct:
\begin{itemize}[leftmargin=1.2em, nosep]
  \item \textbf{stopping at the product} --- $\sqrt6\cdot\sqrt{10}$ left as $\sqrt{60}$; yesterday's ``check the leftover'' failing in a new setting, and the reason the examples were chosen so neither factor simplifies alone;
  \item \textbf{dropping the middle term} --- $\left(\sqrt5+\sqrt2\right)^{2}=7$; do not re-explain, just write $(x+4)^{2}$ next to it and wait;
  \item \textbf{multiplying only the denominator} by the conjugate --- the tell that a student thinks rationalizing is something you do \emph{to the bottom} rather than to the whole fraction;
  \item \textbf{using the same binomial} instead of the conjugate --- the work \emph{looks} right and accomplishes nothing, which is exactly why it is the exit ticket's hardest distractor;
  \item \textbf{the sign of a negative conjugate product} --- $\frac{6}{\sqrt3-3}$ gives $3-9=-6$, and the minus must distribute over \emph{both} numerator terms (homework 3(h));
  \item \textbf{wrong multiplier at a higher index} --- $\frac{1}{\sqrt[3]2}\cdot\frac{\sqrt[3]2}{\sqrt[3]2}$; the same misconception as calling $\sqrt[3]{9}$ unfinished on the 6.2 exit ticket;
  \item \textbf{forgetting to reduce} --- $\frac{6\left(3+\sqrt5\right)}{4}$ left unreduced.
\end{itemize}
Cold-call prompts: ``You have an answer --- is it \emph{simplest}?'' ``How many products does FOIL
give you? Show me all four.'' ``What are you multiplying the whole fraction by, and what does that
number equal?'' ``Is the denominator rational \emph{yet}?'' ``What is the index, and how many
copies do you need under there?''
\end{skillbox}

\begin{skillbox}[Group Work \& Differentiation]{redbox}
\textbf{Making It Rational} activity (tiered; mirrors the notes).
\begin{multicols}{3}
\textbf{Tier R --- Remediate.} A five-row multiply-and-simplify scaffold whose payoff is the
follow-up question --- \emph{three} rows finished with no radical at all, and the reason is that the
radicands multiplied to a perfect power matching the index. Then distributing, a
``what do I multiply by?'' table for monomial denominators (one row, $\frac{5}{\sqrt5}=\sqrt5$,
reduces further than students expect), and a conjugate table ending in the question of why the
product is always radical-free.

\columnbreak
\textbf{Tier A --- Approaching.} Four products including an algebraic one and a cube root; four
expansions where (c) and (d) are deliberately paired so that only one comes out rational; four
rationalizations spanning monomial, simplify-first, and both binomial types; a three-row
\textbf{error analysis} --- the highest-value items on the page, worth a whole-class share; and a
\textbf{rectangle} with sides $3\pm\sqrt5$ whose area (\,$4$\,) \emph{and} perimeter (\,$12$\,) are
both rational for two \emph{different} reasons, closed by a square whose area is not.

\columnbreak
\textbf{Tier E --- Extension.} Three parts. \textbf{(1)} Rationalize $\frac{1}{3-\sqrt5}$ three
ways, two of which fail --- so the conjugate is \emph{ruled in} rather than announced --- then the
general proof and the 3.4 sign question. \textbf{(2)} Rationalizing the \textbf{numerator}:
$\frac{\sqrt{x+9}-3}{x}$ is $\frac00$ at $x=0$, becomes $\frac{1}{\sqrt{x+9}+3}$, and a table at
$x=7,1,0.1,0.01$ closes on $\frac16$ --- a genuine calculus limit, done with today's tool.
\textbf{(3)} Higher-index rationalizing ($\sqrt[3]2$, $\sqrt[3]9$, $\sqrt[4]8$), generalized, then
re-derived with rational exponents.
\end{multicols}
\end{skillbox}

\begin{skillbox}[Individual Work \& Assessment]{redbox}
\textbf{Exit Ticket} (independent, no notes, no calculator): four computations spanning a product
that must be simplified afterward ($\sqrt6\cdot\sqrt{18}=6\sqrt3$), a FOIL, a monomial
rationalization, and a binomial one chosen so the denominator collapses completely
($\frac{2}{\sqrt7-\sqrt5}=\sqrt7+\sqrt5$) --- a wrong answer there is about method, not stamina; an
\textbf{explanation} item requiring the student to show $\frac{1}{\sqrt2}=\frac{\sqrt2}{2}$ and say
why multiplying by $\frac{\sqrt2}{\sqrt2}$ is safe; and an \textbf{SOL-style multiple-choice} item
on $\frac{4}{2-\sqrt3}$ (answer $8+4\sqrt3$) whose distractors each break one idea ---
$\frac{8+4\sqrt3}{7}$ added instead of subtracting in $a^{2}-b^{2}$, $8-4\sqrt3$ used the same
binomial rather than the conjugate, $2+\sqrt3$ cancelled the $4$ against the $2$. Sort into
``got it / did not simplify / middle-term droppers / conjugate confusion.'' The last pile governs
how much of 6.4 opens with algebra; the explanation item, not any computation, is what identifies
students who are performing rationalizing as a ritual and will therefore not reach for it in Lesson
6.5, where nobody prompts for it by name.
\end{skillbox}

\begin{skillbox}[Reinforcement \& Extension]{goldbox}
\textbf{Homework:} eight products (including $\left(3\sqrt7\right)^{2}$, a cube root, an algebraic
radicand, and $2\sqrt{12}\cdot\sqrt{27}=36$); six expansions closed by a diagnostic question ---
two of the six are conjugate pairs and \emph{predictable} in advance, while a third, $\sqrt2
\left(\sqrt8+\sqrt{18}\right)=10$, is rational only by luck of the numbers, and students who cannot
tell those apart have learned a pattern instead of a mechanism; eight divide-or-rationalize items
ending on $\frac{6}{\sqrt3-3}$, whose conjugate product is \emph{negative}; a two-part conceptual
item (refute $\left(\sqrt3+\sqrt5\right)^{2}=8$ numerically; explain why $2-\sqrt3$ works where
$2+\sqrt3$ fails \emph{and} why either is allowed); and a closing ``two roads'' item that
rationalizes $\frac{1}{\sqrt[3]5}$ with radicals and again with rational exponents ---
$5^{-1/3}=\frac{5^{2/3}}{5}$ --- so that rationalizing is revealed as nothing but moving a negative
exponent up. \textbf{Extension:} the \textbf{golden ratio} $\varphi=\frac{1+\sqrt5}{2}$. Unlike the
unit's earlier models (free-fall, Kepler, Kleiber, the pendulum), rationalizing here is not a
tidying-up step --- it is the \emph{only} way to see the result. A table produces the startling
$1.6180 / 0.6180 / 2.6180$; the conjugate then proves $\frac1\varphi=\varphi-1$ and
$\varphi^{2}=\varphi+1$; and the pre-drawn golden rectangle closes it: cut off the $1\times1$
square and the leftover has ratio $\frac{1}{\varphi-1}=\frac{2}{\sqrt5-1}=\varphi$, so the cut can
be repeated forever. \textbf{Preview:} Lesson 6.4, \emph{Graphing Radical Functions \&
Transformations} --- the algebra is finished, and the unit returns to the pictures from 6.0 with
$y=a\sqrt{x-h}+k$ and $y=a\sqrt[3]{x-h}+k$, the endpoint as the transformation anchor, in both
directions (equation $\to$ graph and graph $\to$ equation).
\end{skillbox}
\end{document}
